\documentclass[11pt]{article}
\usepackage{fontspec}
\usepackage[margin=1in]{geometry}
\usepackage{setspace}
\onehalfspacing
\setmainfont{TeX Gyre Termes} % or any system font you prefer

\begin{document}

\begin{center}
    \textbf{\LARGE Statement of Purpose} \\
    \vspace{0.2em}
    \large Matin Badrizade\\ B.Sc. Student in Computer Engineering \\ Amirkabir University of Technology (Tehran Polytechnic)
\end{center}

Ever since I started preparing for Informatics Olympiads, I’ve become increasingly interested in problem solving and in using rigorous mathematical tools and critical thinking to do so. I’ve since continued to apply this mindset across various fields, and I’m fascinated by how new ideas arise every day to solve problems in different settings: mathematics, algorithms, statistics, software engineering, and beyond.

The Informatics Olympiad taught me more than just algorithms and data structures; it instilled a mindset: to break down complex problems, reason abstractly, and construct elegant, provably correct solutions under constraints. This intellectual discipline has shaped not only how I approach challenges in competitive programming, but also how I view the value of critical thinking and rigorous reasoning in real-world contexts.

However, that was not the end of my competitive programming journey. Over the past two years as a student preparing for ICPC contests, I’ve learned many new topics—examples include max-flow algorithms, computational geometry, and FFT-like algorithms—and these continue to amaze me with their elegance.

I’ve also gotten into cryptography and its basic theoretical foundations, encountered the limits of computation, and am more convinced than ever that I want to keep exploring ideas and, hopefully, contribute to future ones.

I recognize that solving predefined problems—however challenging—is only one pillar of intellectual growth. Research represents the next frontier: the opportunity to ask original questions, explore uncharted territory, and contribute—not just apply—knowledge. I am particularly drawn to problems in algorithm design, complexity theory, cryptography, and computer systems generally.

This leads me to strongly believe that the next step for me is to gain research experience, so I can satisfy my passion for problem solving and move forward in my academic journey.
\end{document}